% sec1-intro.tex -- Introduction. Contributions map 1:1 to T1/T2/T3; related work in 1.3.
\section{Introduction}\label{sec:intro}

In online vector balancing, vectors $v_1,\dots,v_T\in\R^n$ arrive one at a time, and each must be assigned a label as soon
as it appears so as to keep the labeled partial sums balanced at every prefix. The two-label (signing) version is the online
form of the Beck--Fiala and Koml\'os problems; here we study the $k$-label (coloring) version, where the imbalance at a
coordinate is the gap between the largest and smallest color load there.

We work in the average-case model of Altschuler and Tikhomirov~\cite{AT2025}: the arrivals are i.i.d.\ uniform $0/1$ vectors
with exactly $d$ ones, $T=\Th(n)$, and $2\le d\le(\loglog n)^2/\log\log\log n$. They proved that the optimal \emph{online}
two-color prefix discrepancy is $\Th(\loglog n)$, independent of $d$, an exponential separation from the offline value
$\Th(\sqrt d)$ of random $d$-regular set systems. We ask how the number of colors $k$ changes this, and answer it in two of
the three regimes of $k$, reducing the third to a single high-probability lemma.

Write $\lam=Td/(nk)=\Th(d/k)$ for the per-color occupancy, $\Hscale=\loglog n$, and, for a scale parameter $\Lscale$,
$\Psifn(\lam,\Lscale)=1+\sqrt{\lam\Lscale}+\Lscale/\log(e+\Lscale/\lam)$; the offline threshold uses $\Lscale=\log k$.

\paragraph{Results.}
Our first result determines the \emph{offline} value up to constant factors, in the growing-palette regime. It is the spread
of $k$ independent Poisson loads at a coordinate, a balls-into-bins occupancy threshold.

\begin{itemize}
  \item[\textbf{(T1)}] (\Cref{thm:offline}) For $C_0\log^5(dk)\le k\le d$ with $d,k=(\log n)^{o(1)}$ and $\lam\ge1$, the
  optimal offline (final) discrepancy is $\offdisc_k=\Th(\Psifn(\lam,\log k))$.
\end{itemize}
The three summands of $\Psifn$ give three regimes: $\sqrt{\lam\log k}$ for $\lam\gtrsim\log k$, then
$\log k/\log((\log k)/\lam)$, then $\log k/\loglog k$ for $\lam=\Th(1)$. Offline, more colors help: the value drops from
$\Th(\sqrt d)$ at $k=2$ toward this much smaller occupancy threshold.

Our second result determines the \emph{online} value for every fixed number of colors.

\begin{itemize}
  \item[\textbf{(T2)}] (\Cref{thm:online}) For every fixed $k\ge2$, $\ondisc_k=\Th(\loglog n)$, independent of $d$.
\end{itemize}
Online, more colors do not help when $k$ is constant: the $\Th(\loglog n)$ barrier of \cite{AT2025} persists for all fixed
$k$, and the upper bound $\Oh(\loglog n)$ in fact holds for all $k$.

The two results give a multicolor online--offline separation.

\begin{itemize}
  \item[\textbf{(T3)}] (\Cref{thm:separation}) For every fixed $k$, $\ondisc_k=\Th(\loglog n)\gg\offdisc_k=\Oh(\sqrt d)=o(\loglog n)$.
\end{itemize}

\begin{table}[t]
\centering
\renewcommand{\arraystretch}{1.3}
\begin{tabular}{l l l l}
\hline
regime of $k$ & offline $\offdisc_k$ & online $\ondisc_k$ & where \\
\hline
$k=2$ & $\Th(\sqrt d)$ & $\Th(\loglog n)$ & \cite{AT2025}, \cite{BansalMeka2019} \\
fixed $k\ge 2$ & $\Oh(\sqrt d)$ & $\Th(\loglog n)$ & \Cref{thm:separation} \\
$C_0\log^5(dk)\le k\le d,\ \lam\ge1$ & $\Th(\Psifn(\lam,\log k))$ & $[\max\{\Psifn(\lam,\log k),\tfrac{\loglog n}{k}\},\,\loglog n]$ & \Cref{thm:offline}, \Cref{cor:bracket} \\
\hline
\end{tabular}
\caption{The offline value collapses to the occupancy threshold $\Psifn(\lam,\log k)$ as $k$ grows, while the online value
stays at the $\Th(\loglog n)$ barrier for every fixed $k$; here $\lam=\Th(d/k)$ and $d,k=(\log n)^{o(1)}$. The growing-$k$
online order is bracketed (\Cref{sec:open}).}\label{tab:results}
\end{table}

\subsection{The growing-$k$ regime}\label{sec:intro-growing}

For growing $k$ the online value is bracketed but not pinned: under the hypotheses of \Cref{cor:bracket}
($C_0\log^5(dk)\le k\le d$, $\lam\ge1$), $\max\{\Psifn(\lam,\log k),\loglog n/k\}\le\ondisc_k\le\loglog n$. The online-specific
term $\loglog n/k$ exceeds the offline floor for small $k$, while the floor $\Psifn$ takes over as $k$ approaches $d$. We reduce the exact order to one high-probability lemma about an asymmetric
online rule, and show that the symmetric approach provably cannot reach it (\Cref{sec:open}). We conjecture
$\ondisc_k=\Th(\max\{\Psifn(\lam,\log k),\loglog n/k\})$.

\subsection{Techniques}\label{sec:intro-techniques}

\paragraph{The offline upper bound.}
The technical heart is producing, from a single draw of the incidence matrix, one $k$-coloring whose every coordinate
range is $\Oh(\Psifn)$. Three ingredients combine (\Cref{sec:offline}). A structural lemma certifies that the matrix is
simultaneously well behaved for \emph{every} column subset, using negative association so that an adaptive recursion never
spends fresh randomness. A density-sensitive splitting lemma halves any column class into prescribed sizes with the
occupancy-scale error, by running the Lovett--Meka partial-coloring walk with locally tuned thresholds and an all-ones
equality that keeps the sizes exact. A recursive color-tree composes these along a balanced binary tree; an error at a node
with $q$ leaves attenuates by $\Oh(1/q)$, and the geometric path sum reassembles exactly into $\Psifn$, with no extra factor
of $\log d$ or $\loglog n$. The matching lower bound is a first-moment Poissonization that reads off the spread of $k$
Poisson loads.

\paragraph{The online upper bound.}
The fixed-$k$ upper bound extends the seed-and-repair scheme of \cite{AT2025} to $k$ colors. The algorithm colors by a random
seed, repairing only the uniquely largest exceptional coordinate of each arrival by a minimum-load color, which never raises
its own range. The category and double-exponential machinery of the two-color analysis uses only independence of future
arrivals, sparsity of the seed-exceptional set, and the implication ``the chosen color differs from the seed only when an
already-exceptional coordinate is present''. None of this refers to the number of colors, so it carries over verbatim, which
we make precise in \Cref{app:online}.

\subsection{Related work}\label{sec:intro-related}

\paragraph{Offline discrepancy.}
The combinatorial discrepancy of a set system asks for a two-coloring keeping every set balanced; the classical bounds are
Beck--Fiala's $2t-1$ for degree $t$ \cite{BeckFiala1981}, Spencer's $\Oh(\sqrt n)$ for $n$ sets \cite{Spencer1985}, the
Bárány--Grinberg vector-balancing bound \cite{BaranyGrinberg1981}, and Banaszczyk's $\Oh(\sqrt{\log n})$ for the Komlós and
Beck--Fiala settings \cite{Banaszczyk1998}; see the surveys and texts \cite{BeckSos1995,Matousek1999,Chazelle2000,AlonSpencer2016}.
A line of constructive algorithms now matches these existentially: Bansal's semidefinite-programming (SDP) rounding \cite{Bansal2010}, the Lovett--Meka
partial-coloring walk \cite{LovettMeka2015}, Rothvoss's convex-set method \cite{Rothvoss2017,EldanSingh2018}, the Gram--Schmidt
walk and its relatives \cite{BansalGarg2017,BansalDadushGargLovett2018}, and an algorithm matching Banaszczyk's Komlós bound
\cite{BansalDadushGarg2019,Nikolov2013}. The structure of the problem is by now well mapped: hereditary discrepancy and its
determinant and $\gamma_2$ lower bounds \cite{LovaszSpencerVesztergombi1986,Matousek2013,MatousekNikolov2015,Larsen2019},
matching hardness \cite{CharikarNewmanNikolov2011}, the refutation of Beck's three-permutations conjecture
\cite{NewmanNeimanNikolov2012}, deterministic constructions \cite{BansalSpencer2013}, and the classical lower-bound technology
\cite{Beck1981,Spencer1987}. For random set systems the picture is sharper than the worst case: the discrepancy of a random
$d$-regular system is $\Th(\sqrt d)$, due to Bansal--Meka, Ezra--Lovett, and Hoberg--Rothvoss
\cite{BansalMeka2019,EzraLovett2016,HobergRothvoss2019}. We use \cite{LovettMeka2015} as our splitting primitive.

\paragraph{Online and stochastic discrepancy.}
Online vector balancing dates to Spencer's balancing games \cite{Spencer1977}. The recent worst-case-adversary and
stochastic theory is extensive: online thinning \cite{DwivediFGGR2019}, the self-balancing walk \cite{AlweissLS2021}, the
Gaussian fixed-point walk \cite{LiuSahSawhney2022}, online geometric discrepancy \cite{BansalJSS2020}, prefix discrepancy and
smoothed analysis \cite{BansalPrefix2022,BansalKomlos2022}, optimal online discrepancy \cite{KulkarniRR2024}, and balancing
sums of random vectors \cite{AruNSV2018}. Online \emph{stochastic} discrepancy for general distributions gives polylogarithmic
prefix bounds \cite{BansalStochastic2021}. Closest to us, Altschuler and Tikhomirov \cite{AT2025} pinpoint the two-color value
in the sparse i.i.d.\ model at $\Th(\loglog n)$, independent of $d$, an exponential separation from the offline $\Th(\sqrt d)$;
our paper is the $k$-color extension of their model.

\paragraph{Multicolor discrepancy and balls-into-bins.}
Multicolor (or $k$-color) discrepancy was introduced by Doerr and Srivastav \cite{DoerrSrivastav2003}; classical reductions to
the two-color case lose factors and do not yield the $1/\sqrt k$ average-case gain of \Cref{thm:offline}. In the worst case the
value is governed by lower bounds of order $\sqrt n$ \cite{MulticolorLB2025}, the opposite dependence on $k$; online multicolor
discrepancy for general unit vectors is $\Th(\sqrt{\log T})$ and is equivalent to online envy minimization \cite{HPVX2025}.
This equivalence ties our question to online fair division, where the goal is to keep allocations envy-free over time
\cite{LiptonMarkakisMosselSaberi2004,BenadeKPP2018,Caragiannis2019}. Our occupancy threshold $\Psifn$ is exactly the
maximum-load law of (weighted) balls-into-bins allocation
\cite{ABKU1999,BCSV2006,RaabSteger1998,TalwarWieder2007,Gonnet1981,KolchinSevastyanovChistyakov1978}, and the growing-$k$
obstruction in \Cref{sec:open} is the same staggered-Fibonacci mechanism behind \emph{Always-Go-Left} and the power of two
choices \cite{Vocking2003,MitzenmacherUpfal2017,Mitzenmacher2001,PeresTalwarWieder2015}.

\paragraph{Tools.}
Beyond \cite{LovettMeka2015}, our offline construction rests on negative association and strong Rayleigh measures
\cite{JoagDevProschan1983,DubhashiRanjan1998,Pemantle2000,BorceaBrandenLiggett2009,BrandenHuh2020}, whose concentration theory
\cite{PemantlePeres2014} and log-concave-polynomial calculus \cite{AnariOveisGharanVinzant2018} underpin the all-subsets
union, and on Bennett-type Poisson tails and the method of bounded differences
\cite{Bennett1962,BoucheronLugosiMassart2013,McDiarmid1989,Janson2004,Vershynin2018,DubhashiPanconesi2009}; the recursive
exact-cardinality rounding parallels online dependent rounding for matchings \cite{NSW2025}.
